\documentclass[11pt,a4paper]{article}
\usepackage[utf8]{inputenc}
\usepackage[french]{babel}
\usepackage[T1]{fontenc}
\usepackage{geometry}
\geometry{top=2.5cm, bottom=2.5cm, left=2.2cm, right=2.2cm}
\usepackage{xcolor}
\usepackage{hyperref}
\usepackage{listings}
\usepackage{tcolorbox}
\tcbuselibrary{skins,breakable}
\usepackage{tikz}
\usetikzlibrary{shapes.geometric, arrows.meta, positioning, shadows, fit}
\usepackage{array}
\usepackage{booktabs}
\usepackage{fancyhdr}
\usepackage{float}

% Palette de couleurs professionnelle
\definecolor{navyprimary}{RGB}{15, 23, 42}
\definecolor{blueaccent}{RGB}{29, 78, 216}
\definecolor{slategray}{RGB}{71, 85, 105}
\definecolor{codebg}{RGB}{248, 250, 252}
\definecolor{codeframe}{RGB}{226, 232, 240}
\definecolor{greenaccent}{RGB}{22, 101, 52}
\definecolor{purplaccent}{RGB}{126, 34, 206}
\definecolor{amberaccent}{RGB}{180, 83, 9}

% Configuration de hyperref
\hypersetup{
    colorlinks=true,
    linkcolor=blueaccent,
    filecolor=purplaccent,
    urlcolor=blueaccent,
    pdftitle={Guide de Présentation Code & Persistance - IncidentFlow},
    pdfauthor={NetMar Technologies},
    pdfpagemode=UseOutlines
}

% Configuration de headheight pour fancyhdr
\setlength{\headheight}{14.5pt}
\pagestyle{fancy}
\fancyhf{}
\rhead{\textcolor{slategray}{\small IncidentFlow --- Reperes Code \& Soutenance}}
\lhead{\textcolor{slategray}{\small Guide de Présentation Oral}}
\rfoot{\textcolor{slategray}{\small Page \thepage}}
\lfoot{\textcolor{slategray}{\small NetMar Technologies}}
\renewcommand{\headrulewidth}{0.4pt}

% Définition de style pour les boîtes colorées
\newtcolorbox{profbox}[1][]{
  colback=blueaccent!5,
  colframe=blueaccent,
  fonttitle=\bfseries\color{white},
  coltitle=white,
  title=Discours à présenter au Professeur / Encadrant,
  enhanced,
  drop shadow,
  #1
}

\newtcolorbox{codebox}[1][]{
  colback=codebg,
  colframe=codeframe,
  fonttitle=\bfseries\color{navyprimary},
  enhanced,
  #1
}

% Listings pour le code
\lstdefinelanguage{CustomJava}{
  language=Java,
  keywordstyle=\color{blueaccent}\bfseries,
  commentstyle=\color{slategray}\itshape,
  stringstyle=\color{greenaccent},
  numberstyle=\tiny\color{slategray},
  breaklines=true,
  showstringspaces=false,
  basicstyle=\ttfamily\small,
  tabsize=2
}

\lstset{
    backgroundcolor=\color{codebg},
    frame=single,
    rulecolor=\color{codeframe},
    captionpos=b,
    numbers=left,
    numbersep=8pt,
    xleftmargin=12pt
}

\begin{document}

% Page de Titre
\begin{titlepage}
    \centering
    \vspace*{1.5cm}
    
    {\Huge \bfseries \color{navyprimary} IncidentFlow}\\[0.4cm]
    {\LARGE \color{blueaccent} Guide de Présentation Code \& Persistance}\\[0.2cm]
    {\large \color{slategray} Repères Fichier par Fichier pour la Soutenance}\\[1.5cm]
    
    \rule{\linewidth}{0.5mm}\\[0.6cm]
    {\Large \bfseries Affichage Frontend (React) \& Persistance Backend (Spring Boot / JPA / PostgreSQL)}\\[0.3cm]
    {\large \textbf{Localisation des Données des Utilisateurs, Rôles \& Workflows Dynamiques}}\\[0.6cm]
    \rule{\linewidth}{0.5mm}\\[2cm]
    
    \begin{minipage}{0.45\textwidth}
        \begin{flushleft} \large
            \textbf{Auteur :}\\
            NetMar Technologies\\
            Équipe IncidentFlow
        \end{flushleft}
    \end{minipage}
    \begin{minipage}{0.45\textwidth}
        \begin{flushright} \large
            \textbf{Destination :}\\
            Démonstration \& Soutenance\\
            Encadrant / Professeur
        \end{flushright}
    \end{minipage}

    \vfill
    {\large Juillet 2026}
\end{titlepage}

\tableofcontents
\newpage

% 1. DISCOURS SYNTHÉTIQUE POUR LE PROFESSEUR
\section{Discours de Présentation Synthétique (À dire à l'oral)}

Lorsque le professeur ou l'encadrant vous demande : \textit{« Montrez-moi dans le code où les données s'affichent et où elles sont persistées »}, vous pouvez présenter l'architecture avec la démonstration suivante :

\begin{profbox}
« Monsieur, l'application est découpée en deux couches principales :

1. \textbf{Côté Frontend (React)} : Les données (Workflows, Utilisateurs, Rôles) sont chargées depuis l'API REST via la méthode \texttt{fetchWorkflows()} dans \texttt{App.jsx}. Elles sont stockées dans des états React (\texttt{useState}). L'utilisateur peut modifier en temps réel le rôle autorisé pour chaque transition dans la liste déroulante (\texttt{<select>}). Lors du clic sur Enregistrer, la fonction \texttt{handleSaveWorkflowGlobally()} transmet le JSON au backend par une requête HTTP \texttt{PUT}.

2. \textbf{Côté Backend (Spring Boot \& PostgreSQL)} : Le contrôleur REST \texttt{WorkflowController} réceptionne la requête sur \texttt{PUT /api/workflows/\{id\}} et la transmet à \texttt{WorkflowService}. Le service exécute une transaction JPA (\texttt{@Transactional}) qui nettoie les anciennes entités enfant (\texttt{orphanRemoval}) et persiste les nouvelles entités \texttt{WorkflowState} et \texttt{WorkflowTransition} (comprenant la colonne \texttt{role\_required}) directement dans PostgreSQL via Spring Data JPA. »
\end{profbox}

---

% 2. COUCHE FRONTEND : AFFICHAGE & ÉDITION (REACT)
\section{Couche Frontend : Affichage et Édition (React UI)}

\subsection{1. Stockage dans l'État React (\texttt{useState}) et Chargement REST}
Les données du workflow actif et de la liste des utilisateurs et rôles sont gérées dans l'état principal du composant React.

\begin{codebox}[title=Localisation dans le Code :]
\textbf{Fichier :} \small\texttt{frontend/src/App.jsx} (Lignes 196 à 199 \& Lignes 825 à 837)
\end{codebox}

\begin{lstlisting}[language=CustomJava, caption={Stockage local et appel API de recuperation dans App.jsx}]
// Lignes 196-199 : Etats React pour les workflows, utilisateurs et roles
const [workflows, setWorkflows] = useState([]);
const [activeWorkflow, setActiveWorkflow] = useState(null);
const [rolesList, setRolesList] = useState([]);

// Lignes 825-837 : Chargement depuis le Backend via API REST
const fetchWorkflows = async () => {
  const res = await fetch(`${API_BASE}/workflows`, { headers: getHeaders() });
  if (res.ok) {
    const data = await res.json();
    setWorkflows(data);
    if (data.length > 0) setActiveWorkflow(data[0]);
  }
};
\end{lstlisting}

\subsection{2. Affichage et Modification Dynamique du Rôle par Transition}
Chaque transition (\texttt{fromState} $\rightarrow$ \texttt{toState}) est affichée avec une liste déroulante interactive permettant d'assigner le rôle autorisé (\texttt{roleRequired}).

\begin{codebox}[title=Localisation dans le Code :]
\textbf{Fichier :} \small\texttt{frontend/src/App.jsx} (Lignes 1457 à 1468 \& Lignes 4236 à 4280)
\end{codebox}

\begin{lstlisting}[language=CustomJava, caption={Mise a jour dynamique de l'etat local dans App.jsx}]
// Lignes 1457-1468 : Fonction de mise a jour du role autorise
const handleUpdateTransitionRole = (fromState, toState, roleRequired) => {
  if (!activeWorkflow) return;
  const updatedTrans = activeWorkflow.transitions.map(t => {
    if (t.fromState === fromState && t.toState === toState) {
      return { ...t, roleRequired: roleRequired || null };
    }
    return t;
  });
  const updatedWf = { ...activeWorkflow, transitions: updatedTrans };
  setActiveWorkflow(updatedWf);
};
\end{lstlisting}

\begin{lstlisting}[language=CustomJava, caption={Rendu JSX de la liste deroulante de role par transition dans App.jsx}]
// Lignes 4236-4280 : Rendu HTML/JSX de la liste des transitions avec <select>
{activeWorkflow.transitions.map((t, idx) => (
  <div key={idx} className="transition-row">
    <span>{t.fromState} -> {t.toState}</span>
    
    <!-- Liste deroulante du Role Autorise -->
    <select
      value={t.roleRequired || ''}
      onChange={(e) => handleUpdateTransitionRole(t.fromState, t.toState, e.target.value)}
    >
      <option value="">Tous les utilisateurs</option>
      <option value="Administrateur">Administrateur</option>
      <option value="Responsable">Responsable</option>
      <option value="Operateur">Operateur</option>
    </select>
  </div>
))}
\end{lstlisting}

\subsection{3. Envoi de la Sauvegarde au Backend}

\begin{codebox}[title=Localisation dans le Code :]
\textbf{Fichier :} \small\texttt{frontend/src/App.jsx} (Lignes 1310 à 1339)
\end{codebox}

\begin{lstlisting}[language=CustomJava, caption={Envoi HTTP PUT vers l'API Backend dans App.jsx}]
// Lignes 1310-1339 : Publication globale du workflow vers le Backend
const handleSaveWorkflowGlobally = async () => {
  const res = await fetch(`${API_BASE}/workflows/${activeWorkflow.id}`, {
    method: 'PUT',
    headers: getHeaders(),
    body: JSON.stringify(activeWorkflow)
  });
  const updated = await res.json();
  setActiveWorkflow(updated);
};
\end{lstlisting}

---

% 3. COUCHE BACKEND : PERSISTANCE JPA & POSTGRESQL
\section{Couche Backend : Persistance JPA, Services et PostgreSQL}

\subsection{1. Point d'Entrée HTTP (Contrôleur REST)}

\begin{codebox}[title=Localisation dans le Code :]
\textbf{Fichier :} \small\texttt{backend/src/main/java/com/netmar/incidentflow/controller/WorkflowController.java} (Lignes 40 à 44)
\end{codebox}

\begin{lstlisting}[language=CustomJava, caption={Endpoint PUT d'enregistrement dans WorkflowController.java}]
@PutMapping("/{id}")
public Workflow updateWorkflow(@PathVariable Long id, @RequestBody Workflow workflowDetails) {
    workflowDetails.setId(id);
    return workflowService.saveWorkflow(workflowDetails);
}
\end{lstlisting}

\subsection{2. Logique Transactionnelle de Persistance (\texttt{WorkflowService})}
Le service exécute l'enregistrement en base de données avec nettoyage des entités orphelines et gestion des versions.

\begin{codebox}[title=Localisation dans le Code :]
\textbf{Fichier :} \small\texttt{backend/src/main/java/com/netmar/incidentflow/service/WorkflowService.java} (Lignes 51 à 160)
\end{codebox}

\begin{lstlisting}[language=CustomJava, caption={Algorithme de persistance transactionnel dans WorkflowService.java}]
@Transactional
@CacheEvict(value = {"workflows", "workflows-category"}, allEntries = true)
public Workflow saveWorkflow(Workflow workflow) {
    validateWorkflowGraph(workflow); // Verification DFS de l'integrite du graphe

    Workflow existing = workflowRepository.findById(workflow.getId())
            .orElseGet(() -> workflow);

    if (existing != workflow) {
        existing.setName(workflow.getName());
        existing.setCategory(workflow.getCategory());
        existing.setActive(workflow.isActive());

        // Supprime les anciens etats et insere les nouveaux (orphanRemoval)
        existing.getStates().clear();
        for (WorkflowState state : workflow.getStates()) {
            existing.getStates().add(WorkflowState.builder()
                    .name(state.getName()).label(state.getLabel())
                    .colorClass(state.getColorClass()).active(state.isActive())
                    .workflow(existing).build());
        }

        // Supprime les anciennes transitions et insere les nouvelles avec roleRequired
        existing.getTransitions().clear();
        for (WorkflowTransition transition : workflow.getTransitions()) {
            existing.getTransitions().add(WorkflowTransition.builder()
                    .fromState(transition.getFromState()).toState(transition.getToState())
                    .roleRequired(transition.getRoleRequired())
                    .requiresComment(transition.isRequiresComment())
                    .workflow(existing).build());
        }
    }
    return workflowRepository.save(existing); // Execution SQL via Hibernate
}
\end{lstlisting}

\subsection{3. Modèle d'Entités JPA (\texttt{Workflow}, \texttt{WorkflowTransition}, \texttt{User}, \texttt{Role})}

\begin{codebox}[title=Localisation dans le Code :]
\textbf{Fichiers :}
\begin{itemize}
    \item \small\texttt{backend/src/main/java/com/netmar/incidentflow/model/Workflow.java} (Lignes 35 à 41)
    \item \small\texttt{backend/src/main/java/com/netmar/incidentflow/model/WorkflowTransition.java} (Lignes 26 à 36)
    \item \small\texttt{backend/src/main/java/com/netmar/incidentflow/model/Role.java} (Lignes 22 à 31)
\end{itemize}
\end{codebox}

\begin{lstlisting}[language=CustomJava, caption={Annotations de cascade JPA dans Workflow.java et WorkflowTransition.java}]
// Dans Workflow.java : Cascade complete et suppression d'orphelins
@OneToMany(mappedBy = "workflow", cascade = CascadeType.ALL, orphanRemoval = true, fetch = FetchType.EAGER)
private List<WorkflowState> states = new ArrayList<>();

@OneToMany(mappedBy = "workflow", cascade = CascadeType.ALL, orphanRemoval = true, fetch = FetchType.EAGER)
private List<WorkflowTransition> transitions = new ArrayList<>();

// Dans WorkflowTransition.java : Propriete du role requis
@Column(name = "role_required")
private String roleRequired;
\end{lstlisting}

---

% 4. TABLEAU RÉCAPITULATIF POUR LA SOUTENANCE
\section{Tableau Récapitulatif des Repères dans le Code}

\begin{table}[H]
\centering
\small
\begin{tabular}{>{\bfseries}p{3.2cm} p{3.8cm} p{5.5cm} p{2.0cm}}
\toprule
\color{navyprimary}Composant & \color{navyprimary}Rôle / Fonction & \color{navyprimary}Fichier & \color{navyprimary}Lignes \\
\midrule
Formulaire UI & Affichage transitions \& rôle & \texttt{frontend/src/App.jsx} & L4236--4280 \\
Mise à jour State & Modification rôle dynamique & \texttt{frontend/src/App.jsx} & L1457--1468 \\
Publication REST & Appel HTTP PUT vers API & \texttt{frontend/src/App.jsx} & L1310--1339 \\
Contrôleur REST & Endpoint \texttt{PUT /api/workflows} & \texttt{WorkflowController.java} & L40--44 \\
Service Persistance & Sauvegarde transactionnelle JPA & \texttt{WorkflowService.java} & L51--160 \\
Entité Workflow & Cascade JPA \& Orphan Removal & \texttt{Workflow.java} & L35--41 \\
Entité Transition & Colonne \texttt{role\_required} & \texttt{WorkflowTransition.java} & L26--36 \\
Matrice RBAC & Jointure Roles/Permissions & \texttt{Role.java} & L22--31 \\
\bottomrule
\end{tabular}
\caption{Tableau de repérage rapide pour les questions du professeur.}
\end{table}

---

% 5. SCHÉMA FLUX DE PERSISTANCE (TIKZ)
\section{Schéma Synthétique du Flux de Persistance}

\begin{figure}[H]
\centering
\begin{tikzpicture}[
    node distance=2cm,
    stepBox/.style={rectangle, draw=navyprimary, fill=codebg, rounded corners, thick, minimum width=4cm, minimum height=1.1cm, align=center, font=\small\bfseries, drop shadow},
    arrow/.style={-{Stealth[scale=1.2]}, thick, draw=blueaccent}
]

\node[stepBox] (step1) {1. Saisie Utilisateur (React)\\ \normalfont\tiny Choix du rôle dans <select>};
\node[stepBox, right=2.5cm of step1] (step2) {2. State React (\texttt{activeWorkflow})\\ \normalfont\tiny Notification du composant local};
\node[stepBox, below=1.8cm of step2] (step3) {3. Requête HTTP PUT\\ \normalfont\tiny Transmet le JSON complet à Spring};
\node[stepBox, left=2.5cm of step3] (step4) {4. \texttt{WorkflowController}\\ \normalfont\tiny Réception et validation REST};
\node[stepBox, below=1.8cm of step4] (step5) {5. \texttt{WorkflowService}\\ \normalfont\tiny \texttt{@Transactional} + \texttt{orphanRemoval}};
\node[stepBox, right=2.5cm of step5] (step6) {6. Base PostgreSQL\\ \normalfont\tiny SQL \texttt{UPDATE} / \texttt{DELETE} / \texttt{INSERT}};

\draw[arrow] (step1) -- (step2);
\draw[arrow] (step2) -- (step3);
\draw[arrow] (step3) -- (step4);
\draw[arrow] (step4) -- (step5);
\draw[arrow] (step5) -- (step6);

\end{tikzpicture}
\caption{Flux complet d'affichage et de persistance des données.}
\end{figure}

\end{document}
